\documentclass[11pt]{article}
\usepackage[margin=0.9in]{geometry}
\usepackage[T1]{fontenc}
\usepackage{lmodern,booktabs,tabularx,amsmath,hyperref,microtype}
\hypersetup{pdftitle={Supplementary Clarifications and Paired-Outcome Accounting},pdfauthor={},hidelinks}
\renewcommand{\thesection}{S\arabic{section}}
\renewcommand{\thepage}{S\arabic{page}}
\setlength{\parindent}{0pt}
\setlength{\parskip}{6pt}
\begin{document}
\begin{center}
{\Large\bfseries Supplementary Clarifications and\\Paired-Outcome Accounting}\par\medskip
When Valid Evidence Goes Unused in Oncology Agents
\end{center}

This addendum explains the submitted study and summarizes existing terminal
outcomes in a post hoc descriptive analysis. We ran no new model generations.
Original scores, primary contrasts, and the status of primary and exploratory
analyses are unchanged. The original bibliography and Appendices A to F are
included in the same PDF.

\section{Study map and endpoint definitions}

The common-prompt study compares valid nonempty inventories under different
supplied contexts. A separate inventory replay measures completeness on saved
search paths. Case reuse and protocol differences limit comparisons across
studies. They cannot be pooled as independent replications or read as successive
improvements to a single intervention.

\begin{table}[h!]
\small
\caption{Study relationships. Case reuse and protocol differences constrain comparisons across rows. Detailed specifications remain in archival Appendices C to E.}
\begin{tabularx}{\linewidth}{@{}p{0.23\linewidth}XX@{}}
\toprule
Study & Cases and interaction & Role in the evidence\\
\midrule
Earlier feedback pilot & Natural and constructed proposals; original retry protocol & Identifier-only versus diagnostic feedback; preceding exploratory evidence.\\
Feedback $\times$ retrieval crossing & 32 co-alteration and 13 sibling-disease positives; four conditions, two repetitions & Supplies the cases reused in both fixed-pool studies below.\\
Initial fixed exposure & Same 45 positives; ten candidates, two eligible alternatives; no search, one revision; separate negative/pass controls & Co-alteration diagnostic 11/64 versus checklist 10/64; benefit criterion fails.\\
Common-prompt context & Same 45 positives and candidates; five conditions, two repetitions; clarified instruction and message assembly & Primary fresh 54/64 versus diagnostic repair 43/64; other contrasts exploratory.\\
Length-matched live retry & Separate 24 co-alteration, 12 sibling-disease, and 12 negative cases; up to two additional searches and resubmissions & Co-alteration diagnostic 30/48 versus checklist 19/48; not an isolated effect of search access.\\
Inventory and replay & 48 cases per arm; recall over 27 positive cases, three repetitions; replay fixes saved search paths & Measures acquired records omitted during finalization.\\
\bottomrule
\end{tabularx}
\end{table}

\textbf{Valid nonempty inventory rate} is the common Boolean endpoint: an answer
contains at least one claim, and every claim passes the schema, visible-ID,
encoded-eligibility, and copied-field checks. \emph{Valid yield} and
\emph{valid recovery} name this endpoint in fresh and post-rejection settings,
respectively. Neither implies a complete inventory. \textbf{Eligible exposure}
records whether at least one eligible candidate became available.
\textbf{Inventory recall} measures the fraction of the eligible reference set
included in the final inventory. The 43/64 valid repair inventories and the
67.2\% retrieved recall belong to different experiments and denominators.

\section{Paired outcomes and exploratory heterogeneity}

We paired fresh and diagnostic-repair outputs by case and repetition index. Each
pair consists of separate generations. It is not a before-and-after measurement
of one trajectory, and the shared repetition index does not imply a shared model
random seed. This descriptive summary supplements the original analysis, which
averages repetitions within cases before variant-cluster resampling.

\begin{table}[h!]
\centering\small
\caption{Post hoc accounting of all 90 fresh/diagnostic-repair pairs from the existing common-prompt study. Invalid answers and schema failures remain distinct. Zero cells are retained.}
\begin{tabular}{@{}llrr@{}}
\toprule
Fresh outcome & Diagnostic repair & Co-alteration & Sibling disease\\
\midrule
Valid inventory & Valid inventory & 41 & 4\\
Valid inventory & Abstention & 13 & 19\\
Abstention & Abstention & 8 & 2\\
Invalid answer & Valid inventory & 2 & 0\\
Schema failure & Abstention & 0 & 1\\
\midrule
Total pairs & & 64 & 26\\
\bottomrule
\end{tabular}
\end{table}

The difference of eleven valid co-alteration inventories comes from thirteen
pairs favoring fresh and two favoring repair. Eight pairs abstain in both
conditions. Thus, not all 21 repair abstentions correspond to a successful fresh
output. Fresh also produces two invalid affirmative answers, both blocked by the
shared checker. Higher valid yield does not establish a uniform safety
advantage. Attempted invalid answers must remain separate from released errors.

Sibling disease shows a larger exploratory difference: fresh 23/26 versus
diagnostic repair 4/26, or +73.1 percentage points (reported variant-cluster
interval [45.5, 93.3]). Nineteen pairs have a valid fresh inventory and a repair
abstention. This stratum includes thirteen cases and eleven variant clusters.
Its contrast is exploratory and unadjusted for multiple comparisons. Cases,
genes, and candidate content also differ between strata; the result does not
isolate a causal effect of disease rejection versus co-alteration rejection.

The primary co-alteration comparison remains +17.2 points [2.8, 33.3], based on
32 cases, 30 variant clusters, two repetitions, and 10,000 paired bootstrap
draws. At the case-mean level, seven cases favor fresh, 24 tie, and one favors
repair. These case-level counts and the repetition-index pairs describe
different units and should not be interchanged.

\section{Generated explanations and submitted actions}

The score depends on the structured action and claims, which can disagree with
the submitted explanation. In an AI-assisted inspection, the authors selected
the following examples from diagnostic-repair abstentions:

\begin{itemize}
\item In co-alteration case CQ001, repetition one, the explanation ends by
recognizing two candidates as valid, while the action remains \texttt{abstain}
and the claim list is empty.
\item In case B032, repetition two, the explanation recognizes eligible
candidates but explicitly cites concern about the preceding rejection when
choosing to abstain.
\item In case N0443, repetition one, the explanation recognizes encoded
eligibility but raises concerns about implicit treatment-population requirements
that are not encoded in the inventory predicate.
\end{itemize}

In these examples, the explanation acknowledges eligibility even though the
action omits the candidates. A response can also question the clinical or
source-level interpretation of a candidate that meets the encoded inventory
rules. Such caution is not necessarily clinically inappropriate.

We selected these examples after inspecting outcomes and have not completed a
qualitative census or estimated category frequencies. The examples do not
establish an internal causal mechanism. The \texttt{reasoning} field contains
generated text; it does not expose hidden reasoning. Disagreement between that
text and the action does not establish that output-field order caused the
disagreement.

A systematic follow-up audit should separate scope misclassification, added
requirements, explicit caution after rejection, and final-explanation/action
disagreement, retaining an unresolved category. There are 43 diagnostic-repair
abstentions across both positive strata (21 co-alteration and 22 sibling
disease). An audit could compare these abstentions with their paired fresh
outputs, or cover all 180 fresh and diagnostic outputs to include successful
submissions. We have not completed that audit or estimated category frequencies.

Existing audits have different targets. Archival Section C reports a 237-item
enriched AI-reviewed sample from initial fixed exposure and earlier
pilot/crossed explanation audits. The 211 live-retry items prepared for review have not been rated.
Section E.5 reports an inventory-task qualitative sample and a post hoc census
of ten failed inventories. Those studies do not constitute an audit of the
common-prompt 21 co-alteration abstentions.

\section{Prompt sensitivity and identification limits}

Diagnostic recovery rises from 11/64 in initial fixed exposure to 43/64 in the
common-prompt study on the same co-alteration cases and candidate payloads.
Three protocol changes accompany this 50.0-point descriptive difference:

\begin{enumerate}
\item The initial proposal appears in both the user payload and assistant call
in initial fixed exposure, but only once in the later protocol.
\item The common instruction explicitly permits replacement of a rejected record.
\item It explicitly includes eligible prognostic and N/A-therapy records.
\end{enumerate}

Exact initial and common instructions are disclosed in archival Sections C.4 and
D.1. Because the protocol changes were bundled, the comparison cannot attribute
the difference to an individual edit or quantify how much of the earlier failure
came from underspecification. The recovery rate nevertheless differs
substantially between the two formulations. The later fresh versus repair
comparison holds the common instruction fixed and addresses a different
question.

Fresh receives no erroneous proposal or rejection feedback. Repair receives both
as supplied assistant/tool history. This comparison changes the proposal,
feedback, message presentation, and input length as a package. Checklist and
diagnostic \emph{feedback strings} are length-matched within the specified
tolerance; fresh and repair \emph{total contexts} are not length-matched. Reset
retains proposal and feedback content as user review material. The diagnostic
reset and repair totals are both 43/64, but two case means improve, 29 tie, and
one worsens. Equal aggregate yields do not establish equivalence, identical
behavior, or validity for naturally occurring errors.

Separating the components requires new controls, such as matched-length neutral
context and proposal/feedback presence varied within a fixed presentation.
Restarting after removing the rejected ID is a distinct operational policy:
removing one of ten candidates also changes candidate count and input length. A
matched removal of an unrelated ineligible candidate would help distinguish
rejected-ID removal from generic distractor reduction. We have not run these
controls, tested continuations after naturally occurring errors, or replicated
the study across model versions in this update.

We froze the common-prompt analysis locally before generating its outputs, after
seeing the earlier fixed-exposure results. We reused cases and payloads.
Prespecification therefore concerns the subsequent execution and primary
contrast, not outcome-independent selection of the research question or an
independent case sample. The model is one unversioned alias. Shared publications
connect all 45 positives into one component; variant-cluster intervals do not
separately account for that dependence. Their coverage for generalization to new
publication families is not established. A single connected component cannot
supply an informative independent-component bootstrap. Additional repetitions on
the same cases do not resolve these limitations.

\section{Deterministic preservation and the interpretation boundary}

Let $R$ be the deduplicated records acquired on a saved search path, $P$ the
encoded eligibility predicate, and $G$ the full eligible reference set under
that predicate. Serialization emits \[ S=\{r\in R:P(r)\}=R\cap G. \] When
identifier handling and field copying conform to the frozen contract, serialized
recall therefore equals retrieved recall by construction. The empirical result
is the gap between the available eligible pool and the LLM's final inventory;
the identity itself follows from the serialization rule.

On iterative scope-first paths, the LLM retains 847 of 975 available eligible
record exposures and omits at least one eligible record in 29/81 positive
trajectories. The 128 omitted exposures count unique identifiers within each
trajectory. An identifier can recur across trajectories, so this count does not
represent 128 distinct corpus records. Case-macro retrieved and final inventory
recall are 67.216\% and 60.302\%, respectively. Their paired difference is 6.914
points, whereas $847/975$ is pooled retention and uses a different denominator.
Original generation has a 4,096-token response cap with no observed truncations;
serialization is uncapped. The replay fixes saved retrieval and changes
finalization under these output constraints.

When inventory membership is fully computable, deterministic selection and
copying can preserve acquired records without an LLM. This procedure provides a
direct baseline for finalization. Search formulation and source-level
interpretation are possible LLM roles, but this replay does not establish their
necessity, superiority, or clinical benefit. Full-corpus enumeration reaches
100\% encoded recall by construction, without establishing publication-level
validity.

The checker and reference labels use the same frozen curation. Agreement
therefore establishes conformance to the encoded task, but independent
publication support still requires source review. W04 illustrates this boundary
through observations tied to different genotypes and experimental models. It is
one example within an exploratory twelve-family review, not an estimate of
corpus error or validation of the repair effect. Two medical-student readers
used substantive AI assistance and nonblinded joint adjudication; the separate
thirty-pair independent review was discontinued without completed judgments. No
workflow study measured review time or decision quality. Preservation and source
interpretation remain separate evaluation targets.

\clearpage
\section{Finding material cited in the paper}

This guide refers to the 20-page submitted paper. Its condensed appendices and
this supplement's expanded appendices use different subsection and table
numbers. All page numbers below are PDF viewer pages in the combined
\emph{Extended Methods and Results}. On PDF pages 2 to 32, subtract eight from
the printed page number. Tables in addendum sections S1 and S2 have their own
numbering.

\subsection*{Sections and tables}
{\small
\begin{tabularx}{\linewidth}{@{}p{0.38\linewidth}Xr@{}}
\toprule
Reference in the paper & Expanded supplement & PDF pages\\
\midrule
Table 1: protocol roles and budgets & S1 overview; C, Table 14; E.1 & 33, 13, 25\\
B, Table 3: samples and limits & B.1, B.5; study settings in C to E & 7, 10 to 12\\
Introduction; A, Table 2: typed scope audit & A.4, Table 4 & 4 to 5\\
Section 2; B: shared methods & B.1 to B.5 & 7 to 12\\
Section 3; C.1, Table 4: fixed exposure & C.1, Table 15; outcomes in Table 17 & 13 to 15\\
Section 4.3; C.2, Table 4: live retry & C.2, Tables 16 and 17 & 14 to 15\\
C.3: initial prompts and messages & C.4 & 19 to 20\\
Section 4; D.1, Table 5: common prompt and conditions & D.1, Table 22 & 20 to 21\\
D.2: design, scoring and statistics & D.2 & 21 to 22\\
D.3, Tables 6 and 7: outcomes and contrasts & D.3, Tables 23 and 24 & 22 to 23\\
D: detailed resource accounting & D.4, Table 26 & 24\\
Section 5.1; E.1: inventory task & E.1 to E.2, Tables 27 to 29 & 25 to 26\\
E.2, Table 8: saved-pool serialization & E.4, Tables 31 and 32; replay checks & 28 to 29\\
E.3: controls and omissions & E.3 to E.5 & 27 to 29\\
F, Table 9: joint judgments & F.3, Table 35, columns J & 31\\
F, Table 10: W04 observations & F.3, W04 discussion & 31\\
F: source review and its limits & F.1 to F.4, Tables 33 to 36 & 29 to 32\\
\bottomrule
\end{tabularx}
}

The initial system instruction, checklist core, and feedback tail appear in C.4
(PDF pages 19 to 20). The exact common system instruction and final command
appear in D.1 (PDF page 21). These passages let readers compare the disclosed
instructions even though the original prompt JSON files are not included.
Related work is in B.6 (PDF page 12), and the bibliography is on PDF page 2.

\subsection*{What the public data reproduce}
\texttt{data/paired\_outcomes.csv} contains 90 fresh/diagnostic-repair pairs
(180 outputs): 64 co-alteration pairs and 26 sibling-disease pairs.
\texttt{source/analyze.py} recomputes the S2 counts. The CSV covers two of the
five common-prompt conditions, whose full study comprised 450 outputs.
It omits variant-cluster assignments, the other conditions, candidate records,
and full responses. It therefore does not support independent response scoring,
reconstruction of all five conditions, or reproduction of the reported
variant-cluster bootstrap intervals. Original scoring and statistical methods
are described in D.2 to D.3.

\clearpage
\subsection*{Files and documents cited in the paper}

The paths in the paper and original Appendix B.5 refer to the research workspace
used for the experiments. Original Table 13 is a map of that workspace, not this
repository. The original experiment directories, code, and execution records are
not included here. The passages below explain the methods. Running the original experiments would
require the original code and inputs.

{\small
\begin{tabularx}{\linewidth}{@{}p{0.47\linewidth}X@{}}
\toprule
Cited file or document & Description available in this PDF\\
\midrule
\nolinkurl{versions/v14-full-paper-revision/README.md};
\nolinkurl{src/common.py}; \nolinkurl{src/agents.py}
& B.1 to B.5, pages 7 to 12: corpus, eligibility, scoring and execution. Original README and code are not included.\\[4pt]
\nolinkurl{analysis_inputs/prompt_material_exact.json}
& C.4, pages 19 to 20: disclosed initial instruction and feedback text. Original JSON is not included.\\[4pt]
V40 \nolinkurl{src/study.py}; V38 E1 \nolinkurl{src/runner.py} and \nolinkurl{src/feedback.py}
& C.1 to C.2 and C.4, pages 13 to 15 and 19 to 20: repair protocols and message assembly. Original code is not included.\\[4pt]
\nolinkurl{analysis_inputs/v41_public_prompt_contract.json}
& D.1, pages 20 to 21: exact common system instruction and final command. The full JSON, schema and request hashes are not included as files.\\[4pt]
\nolinkurl{versions/v41-repair-history-control/}
& D.1 to D.4, pages 20 to 24: study design and aggregate results. Original experiment directory is not included.\\[4pt]
Other directories, execution plans, timing files and per-run scores listed in B.5
& B.5, pages 10 to 12, and the study appendices. The original files and complete per-run results are not included.\\[4pt]
Saved-pool replay code and inputs listed in B.5
& E.4, pages 28 to 29: procedure and reported verification. Executable replay and original inputs are not included.\\[4pt]
Companion technical report
& A, C and F describe related studies and review scope. The separate report and its additional results are not included.\\[4pt]
Signed reader submissions and source-bearing review materials
& F.1 to F.4, pages 29 to 32: methods and reported judgments. Underlying restricted materials are not included.\\
\bottomrule
\end{tabularx}
}

\subsection*{Building the public PDFs}
\texttt{source/build.py} compiles the cover and addendum, combines them with the
original bibliography and appendices, and verifies that the original pages have
identical text and rendering. It clears identifying PDF metadata and records
checks and artifact hashes in \texttt{data/verification.json}. The original
supplement is retained in \texttt{archive/}. The build requires Python 3,
PyMuPDF, and \texttt{pdflatex}; it does not require the unavailable experiment
files, licensed candidate text, full transcripts, or private reader forms.
The README lists the commands and links to all public data and source files.
\end{document}
